% GENERATED FILE. Edit canonical lessons, then run npm run generate:books.
% canonical-source-hash: fnv1a64:6d0c808b4a35e62a
% canonical-lessons: MR-C369-pradhyapak, MR-C369-abhyasakram, MR-C369-sisyavrtti, MR-C369-sulka, MR-C369-vasatigrh

\chapter{At University}
\label{ch:mr-at-university}
\hlchaptermodality{369}

\begin{chapteropening}
\textbf{By the end of this chapter:} \emph{I can talk about professors, syllabuses, fees and hostels.}

Complete the last lesson of chapter 369: I can talk about professors, syllabuses, fees and hostels.
\end{chapteropening}

\section[\texorpdfstring{\mr{प्राध्यापक} (prādhyāpak)}{प्राध्यापक (prādhyāpak)}]{\mr{प्राध्यापक} (prādhyāpak) — a professor}

\label{lesson:MR-C369-pradhyapak}



\begin{quote}
Before the new one: say the Marathi for to entrust, then the Marathi for to handle.
\end{quote}

\subsection*{You'll want to know: \mr{प्राध्यापक}}
\textbf{\mr{प्राध्यापक}} — \emph{prādhyāpak} — \textquotedblleft{}a professor\textquotedblright{}.

\begin{quote}
\mr{आमचे} \mr{प्राध्यापक} \mr{खूप} \mr{छान} \mr{शिकवतात}.
\end{quote}

\begin{quote}
\emph{āmce prādhyāpak khūp chān śikavtāt.}
\end{quote}

\begin{quote}
Our professor teaches very well.
\end{quote}

Masculine.

Respect takes the plural: \mr{प्राध्यापक} \mr{शिकवतात}.

\begin{cousinweb}
\textbf{\mr{प्राध्यापक}} (\emph{prādhyāpak}) is a learned word built from Sanskrit \emph{pra-}, \textquotedblleft{}forth, before\textquotedblright{}, and \emph{adhyāpaka}, \textquotedblleft{}teacher\textquotedblright{} (CDIAL 278).
\end{cousinweb}

\begin{grammarlens}[title={What you've built}]
One more word for giving reasons and opinions.
\end{grammarlens}

\subsection*{Your turn}
\begin{itemize}
\raggedright
  \item \emph{Say it:} \emph{prādhyāpak}
  \item \emph{Say it:} \emph{prādhyāpak}, once more
  \item \emph{Say it:} \emph{āmce prādhyāpak khūp chān śikavtāt.} — Our professor teaches very well.
  \item \emph{Recall:} say the Marathi for to entrust, then the Marathi for to handle, then say \emph{prādhyāpak} again
\end{itemize}

\begin{tcolorbox}[breakable,colback=teal!4,colframe=teal!35!black,title={Before you move on}]
What does \mr{प्राध्यापक} mean? (\textquotedblleft{}A professor\textquotedblright{}.) Say it once more.
\end{tcolorbox}
\section[\texorpdfstring{\mr{अभ्यासक्रम} (abhyāsakram)}{अभ्यासक्रम (abhyāsakram)}]{\mr{अभ्यासक्रम} (abhyāsakram) — a syllabus; a course of study}

\label{lesson:MR-C369-abhyasakram}



\begin{quote}
Before the new one: say the Marathi for to handle, then the Marathi for a professor.
\end{quote}

\subsection*{You'll want to know: \mr{अभ्यासक्रम}}
\textbf{\mr{अभ्यासक्रम}} — \emph{abhyāsakram} — \textquotedblleft{}a syllabus; a course of study\textquotedblright{}.

\begin{quote}
\mr{या} \mr{वर्षी} \mr{अभ्यासक्रम} \mr{बदलला} \mr{आहे}.
\end{quote}

\begin{quote}
\emph{yā varṣī abhyāsakram badallā āhe.}
\end{quote}

\begin{quote}
The syllabus has changed this year.
\end{quote}

Masculine.

Built on \mr{अभ्यास}, study: the order of study.

\begin{cousinweb}
\textbf{\mr{अभ्यासक्रम}} (\emph{abhyāsakram}) is a learned compound built from Sanskrit \emph{abhyāsa}, \textquotedblleft{}repetition, study\textquotedblright{}, and \emph{krama}, \textquotedblleft{}step, series\textquotedblright{} (Monier-Williams).
\end{cousinweb}

\begin{grammarlens}[title={What you've built}]
One more word for giving reasons and opinions.
\end{grammarlens}

\subsection*{Your turn}
\begin{itemize}
\raggedright
  \item \emph{Say it:} \emph{abhyāsakram}
  \item \emph{Say it:} \emph{abhyāsakram}, once more
  \item \emph{Say it:} \emph{yā varṣī abhyāsakram badallā āhe.} — The syllabus has changed this year.
  \item \emph{Recall:} say the Marathi for to handle, then the Marathi for a professor, then say \emph{abhyāsakram} again
\end{itemize}

\begin{tcolorbox}[breakable,colback=teal!4,colframe=teal!35!black,title={Before you move on}]
What does \mr{अभ्यासक्रम} mean? (\textquotedblleft{}A syllabus; a course of study\textquotedblright{}.) Say it once more.
\end{tcolorbox}
\section[\texorpdfstring{\mr{शिष्यवृत्ती} (śiṣyavṛttī)}{शिष्यवृत्ती (śiṣyavṛttī)}]{\mr{शिष्यवृत्ती} (śiṣyavṛttī) — a scholarship}

\label{lesson:MR-C369-sisyavrtti}



\begin{quote}
Before the new one: say the Marathi for a professor, then the Marathi for a syllabus.
\end{quote}

\subsection*{You'll want to know: \mr{शिष्यवृत्ती}}
\textbf{\mr{शिष्यवृत्ती}} — \emph{śiṣyavṛttī} — \textquotedblleft{}a scholarship\textquotedblright{}.

\begin{quote}
\mr{हुशार} \mr{विद्यार्थ्यांना} \mr{शिष्यवृत्ती} \mr{मिळते}.
\end{quote}

\begin{quote}
\emph{huśār vidyārthyānnā śiṣyavṛttī miḷte.}
\end{quote}

\begin{quote}
Bright students get a scholarship.
\end{quote}

Feminine.

\mr{शिष्य} + \mr{वृत्ती}, pupil + livelihood.

\begin{cousinweb}
\textbf{\mr{शिष्यवृत्ती}} (\emph{śiṣyavṛttī}) is a learned compound built from Sanskrit \emph{śiṣya}, \textquotedblleft{}pupil\textquotedblright{}, and \emph{vṛtti}, \textquotedblleft{}way of life, livelihood\textquotedblright{} (Monier-Williams): a pupil's living.
\end{cousinweb}

\begin{grammarlens}[title={What you've built}]
One more word for giving reasons and opinions.
\end{grammarlens}

\subsection*{Your turn}
\begin{itemize}
\raggedright
  \item \emph{Say it:} \emph{śiṣyavṛttī}
  \item \emph{Say it:} \emph{śiṣyavṛttī}, once more
  \item \emph{Say it:} \emph{huśār vidyārthyānnā śiṣyavṛttī miḷte.} — Bright students get a scholarship.
  \item \emph{Recall:} say the Marathi for a professor, then the Marathi for a syllabus, then say \emph{śiṣyavṛttī} again
\end{itemize}

\begin{tcolorbox}[breakable,colback=teal!4,colframe=teal!35!black,title={Before you move on}]
What does \mr{शिष्यवृत्ती} mean? (\textquotedblleft{}A scholarship\textquotedblright{}.) Say it once more.
\end{tcolorbox}
\section[\texorpdfstring{\mr{शुल्क} (śulka)}{शुल्क (śulka)}]{\mr{शुल्क} (śulka) — a fee}

\label{lesson:MR-C369-sulka}



\begin{quote}
Before the new one: say the Marathi for a syllabus, then the Marathi for a scholarship.
\end{quote}

\subsection*{You'll want to know: \mr{शुल्क}}
\textbf{\mr{शुल्क}} — \emph{śulka} — \textquotedblleft{}a fee\textquotedblright{}.

\begin{quote}
\mr{महाविद्यालयाचं} \mr{शुल्क} \mr{दरवर्षी} \mr{वाढतं}.
\end{quote}

\begin{quote}
\emph{mahāvidyālyācaṁ śulka darvarṣī vāḍhtaṁ.}
\end{quote}

\begin{quote}
The college fee goes up every year.
\end{quote}

Neuter.

Formal; in speech \mr{फी}.

\begin{cousinweb}
\textbf{\mr{शुल्क}} (\emph{śulka}) is a tatsama, taken from Sanskrit \emph{śulka}, \textquotedblleft{}price, value; tax, toll\textquotedblright{} (CDIAL 12543).
\end{cousinweb}

\begin{grammarlens}[title={What you've built}]
One more word for giving reasons and opinions.
\end{grammarlens}

\subsection*{Your turn}
\begin{itemize}
\raggedright
  \item \emph{Say it:} \emph{śulka}
  \item \emph{Say it:} \emph{śulka}, once more
  \item \emph{Say it:} \emph{mahāvidyālyācaṁ śulka darvarṣī vāḍhtaṁ.} — The college fee goes up every year.
  \item \emph{Recall:} say the Marathi for a syllabus, then the Marathi for a scholarship, then say \emph{śulka} again
\end{itemize}

\begin{tcolorbox}[breakable,colback=teal!4,colframe=teal!35!black,title={Before you move on}]
What does \mr{शुल्क} mean? (\textquotedblleft{}A fee\textquotedblright{}.) Say it once more.
\end{tcolorbox}
\section[\texorpdfstring{\mr{वसतिगृह} (vasatigṛh)}{वसतिगृह (vasatigṛh)}]{\mr{वसतिगृह} (vasatigṛh) — a hostel}

\label{lesson:MR-C369-vasatigrh}



\begin{quote}
Before the new one: say the Marathi for a scholarship, then the Marathi for a fee.
\end{quote}

\subsection*{You'll want to know: \mr{वसतिगृह}}
\textbf{\mr{वसतिगृह}} — \emph{vasatigṛh} — \textquotedblleft{}a hostel\textquotedblright{}.

\begin{quote}
\mr{मी} \mr{पहिल्या} \mr{वर्षी} \mr{वसतिगृहात} \mr{राहिलो}.
\end{quote}

\begin{quote}
\emph{mī pahilyā varṣī vasatigṛhāt rāhilo.}
\end{quote}

\begin{quote}
I lived in the hostel in my first year.
\end{quote}

Neuter.

\mr{वसती} + \mr{गृह}, dwelling + house.

\begin{cousinweb}
\textbf{\mr{वसतिगृह}} (\emph{vasatigṛh}) is a learned compound built from Sanskrit \emph{vasati}, \textquotedblleft{}dwelling, staying\textquotedblright{}, and \emph{gṛha}, \textquotedblleft{}house\textquotedblright{} (Monier-Williams).
\end{cousinweb}

\begin{grammarlens}[title={What you've built}]
One more word for giving reasons and opinions.
\end{grammarlens}

\subsection*{Your turn}
\begin{itemize}
\raggedright
  \item \emph{Say it:} \emph{vasatigṛh}
  \item \emph{Say it:} \emph{vasatigṛh}, once more
  \item \emph{Say it:} \emph{mī pahilyā varṣī vasatigṛhāt rāhilo.} — I lived in the hostel in my first year.
  \item \emph{Recall:} say the Marathi for a scholarship, then the Marathi for a fee, then say \emph{vasatigṛh} again
\end{itemize}

\begin{tcolorbox}[breakable,colback=teal!4,colframe=teal!35!black,title={Before you move on}]
What does \mr{वसतिगृह} mean? (\textquotedblleft{}A hostel\textquotedblright{}.) Say it once more.
\end{tcolorbox}
